\section{Motivation}\label{motivation-agentic}

Modern software development is often bottlenecked by the manual cycle of ``hypothesis \(\rightarrow\) implementation \(\rightarrow\) testing''. Agentic Engineering seeks to automate this loop by using LLM-based agents that can write code, run tests, and analyze results autonomously. This approach draws from Karpathy's autoresearch concept~\cite{karpathy2026autoresearch}, which frames AI-driven code modification as an iterative keep/discard loop, and from Optuna's define-by-run hyperparameter optimisation paradigm~\cite{akiba2019optuna}.

\section{Multi-Agent Architecture}\label{multi-agent-architecture}

The project implemented the \textbf{Autolab} system (documented in the project monorepo under \texttt{.opencode/agents/}), a five-agent orchestration layer inspired by Karpathy's autoresearch concept~\cite{karpathy2026autoresearch}. Scientists interact with one primary agent, while specialized subagents handle literature review, queue planning, experiment execution, and reporting.

\subsection{Delegation Map}\label{delegation-map}

\begin{center}
\small
\begin{tabular}{@{}p{2.7cm} p{9.8cm}@{}}
\toprule
\textbf{Agent} & \textbf{Responsibility} \\
\midrule
\textbf{Autolab} & Primary orchestrator --- scientists interact only with this agent, which delegates work to specialized subagents \\
\textbf{Researcher} & Literature scout: searches web/PDFs, extracts methodology, writes hypotheses to \texttt{paper-ideas.md} \\
\textbf{Planner} & Queue coordinator: reads \texttt{paper-ideas.md}, deduplicates hypotheses, assigns run\_ids and branches \\
\textbf{Worker} & Isolated experiment runner: creates git worktree, implements hypothesis, runs \texttt{prepare.py}, records CSV \\
\textbf{Reporter} & Status summarizer: reads CSV and \texttt{results.tsv}, emits F1 summary and staircase report \\
\bottomrule
\end{tabular}
\normalsize
\end{center}

\paragraph{Cost-aware model routing.}
Rather than binding each role to one fixed model vendor, the system routes tasks through \texttt{9router} so each agent can use a model tier that matches the required reasoning strength and cost. In practice, this allows mixing providers such as OpenAI, DeepSeek, OpenCode, MiniMax, NVIDIA, and OpenRouter behind one workflow boundary. ``BigBrain'', ``MidBrain'', and ``MiniBrain'' therefore denote capability tiers rather than one specific model. Example BigBrain choices include GPT-5.4, DeepSeek V4 Pro, and Claude Sonnet 4.8; MidBrain can use models such as DeepSeek V4 Flash, MiniMax M3, or GPT-5.3 Low; MiniBrain can use free or near-free models with weaker reasoning when tasks are repetitive and tightly bounded.

\begin{figure}[H]
\centering
\includegraphics[width=0.92\textwidth,keepaspectratio]{figures/ch04_autolab_agents.png}
\caption{Autolab five-agent interaction: scientist $\rightarrow$ Autolab orchestrator $\rightarrow$ Researcher, Planner, Worker, and Reporter. Worker execution forms a repeatable spawn loop so the planner can launch as many isolated worktree runs as needed. Workers use \texttt{fcntl.flock} for CSV-safe concurrent writes.}
\label{fig:autolab-agents}
\end{figure}

\paragraph{Why experiment-first iteration was preferred.}
A notable shift from earlier research style is that less time was spent manually inspecting special visual traits of the dataset before building methods. Detailed dataset interpretation often depends on slow human feedback, while AI agents can run and compare many experiments quickly at low cost. As a result, the workflow emphasizes early empirical exploration: try more candidate methods sooner, observe which configurations produce strong results, then investigate or operationalize the most effective methods directly. This changes the research loop from ``study first, implement later'' toward ``test early, then deepen where evidence is strongest.''

\paragraph{Why isolated subagents matter.}
Subagents operate in separate git worktrees so parallel tasks do not overwrite the same files. This isolation also supports human-in-the-loop review: after a task completes, its branch can be reviewed or promoted independently, and pull requests can be created without blocking other experiment branches.
\section{Reducing Hallucinations}\label{reducing-hallucinations}

To improve reliability, the workflow uses three guardrails. First, \textbf{Spec-Driven Development (SDD)} uses \texttt{Speckit}~\cite{speckit} together with \texttt{OpenSpec}~\cite{openspec} so agents work against explicit specifications and drift can be detected systematically. Second, \textbf{Test-Driven Development (TDD)} requires an executable check around the change, reducing silent regressions. Third, \textbf{specialized skills} provide bounded instructions for recurring tasks, so the agent follows validated local procedures instead of relying only on broad model memory.

Reusable skills become especially valuable when a workflow repeats and the model otherwise wastes tokens by rediscovering it through trial and error. For example, obtaining a paper PDF may involve searching the web, attempting direct download links, failing on redirects, locating the actual download button, then extracting and reading the final document. Once this pattern was observed repeatedly, encoding it as a reusable skill reduced wasted attempts, token cost, and variance in execution quality.
\section{Overcoming Context Window Limits}\label{overcoming-context-window-limits}

As codebase size and task count increased, controlling token usage and context scope became critical.

\subsection{Reduce-Token-Usage Tools}\label{reduce-token-usage-tools}

The workflow uses explicit token-reduction tools rather than relying only on shorter prompts. One example is Caveman mode, a compressed communication style that removes linguistic overhead while preserving technical content. Another is RTK, used at the proxy layer through \texttt{9router} to reduce prompt and context cost further. These tools are practical engineering responses to long-running multi-agent sessions where unnecessary tokens directly increase latency and monetary cost.

\subsection{Subagents and Worktrees}\label{subagents-and-worktrees}

Isolated git worktrees and a routing layer ensure each subagent receives only task-relevant context. This reduces prompt size, avoids concurrent file-edit conflicts, and lets multiple tasks proceed in parallel without one agent overwriting another's work.

\subsection{Reusable Skills}\label{reusable-skills}

Encoding complex workflows (for example, frontend best practices or FastAPI-specific conventions) into bounded skill definitions constrains behavior to validated patterns and shortens future prompts.

\paragraph{Reference implementations.}
Representative public references for token-reduction and routing ideas include Caveman~\cite{caveman}, 9router~\cite{router9}, and RTK~\cite{rtk}.

\section{Summary}\label{agentic-summary}

Agentic engineering in this project was not treated as a general chatbot wrapper. It was organized as a cost-aware, multi-agent software workflow with explicit routing, isolated execution, reusable skills, and verification guardrails. The result is a research process that can explore more configurations earlier, while still preserving human oversight at branch, review, and promotion boundaries.
